\documentclass[10pt]{article}
\usepackage[margin=2cm]{geometry}
\usepackage[utf8]{inputenc}
\usepackage{cite}
\usepackage{amsmath,amssymb,amsfonts}
\usepackage{graphicx}
\usepackage{textcomp}
\usepackage{booktabs}
\usepackage{array}
\usepackage{longtable}
\usepackage{ragged2e}
\usepackage{placeins}
\usepackage{url}
\renewcommand{\thesection}{S\arabic{section}}
\renewcommand{\thetable}{S\arabic{table}}
\renewcommand{\thefigure}{S\arabic{figure}}
\title{Supplementary Material\\\large Failure-Aware Container Orchestration under Node and Network Partitions in Resource-Constrained Healthcare Cloud Environments}
\author{Kawuma Daniel and Fred Wabendo}
\date{}
\begin{document}
\maketitle
\noindent This supplement holds the four comparative matrices of the literature review (Section~S1, Tables~S1--S5) and the run-level OpenMRS tables summarised in the main paper (Section~S2, Tables~S6--S12). Reference numbers match the main paper. Every run-level value is also in the run directories and CSV files of the repository release.

\section{Comparative Systems Analysis}

\paragraph{Primary search string.} Run in IEEE Xplore, Google Scholar and the ACM Digital Library on 22--23 August 2026; all 20 queries are in the search log.
\begin{quote}
\footnotesize (``container orchestration'' OR ``container orchestrator'' OR Kubernetes OR K8s OR ``container scheduling'' OR ``containerized workload*'' OR ``cloud orchestration'' OR ``edge orchestration'') AND (``node failure'' OR ``host failure'' OR ``server failure'' OR ``node crash'' OR ``network partition'' OR ``network failure'' OR ``connectivity failure'' OR ``link failure'' OR ``packet loss'' OR ``resource exhaustion'' OR ``resource saturation'' OR ``resource contention'' OR ``resource pressure'' OR ``memory exhaustion'' OR ``memory pressure'' OR ``out of memory'' OR ``CPU saturation'' OR ``CPU contention'' OR ``failure detection'' OR ``failure recovery'' OR ``fault tolerance'' OR ``failure-aware'' OR resilien* ) AND (``resource constrained'' OR ``resource-constrained'' OR ``constrained infrastructure'' OR ``limited resources'' OR ``distributed cloud'' OR ``multi-site'' OR ``geo-distributed'' OR ``geographically distributed'' OR ``intermittent connectivity'' OR ``disconnected operation'' OR ``offline operation'' OR ``wide area network'' OR WAN)
\end{quote}
\label{sec:matrices}

This section presents four comparative matrices covering all 24 included papers. In Matrices~1--3 the rows below the rule are the three papers whose decision component is learned \cite{dsouza2026selfhealing, nigam2026chaos, mishra2026aidriven}; they are coded on the same dimensions as the others, and Section~III of the main paper explains how they inform the design. Cross-paper trade-off discussion follows each matrix.

\medskip\noindent\textit{Matrix 1 -- Multi-Dimensional Taxonomy Mapping.}\par\smallskip

The comparative matrices use the following terms after introducing them here in full: wide-area network (WAN), Domain Name System (DNS), service-level agreement (SLA), machine learning (ML), convolutional neural network (CNN), long short-term memory (LSTM), graphics processing unit (GPU), multi-instance GPU (MIG), genetic algorithm (GA), peer-to-peer (P2P), distributed hash table (DHT), mean time to detect (MTTD), mean time to recover (MTTR), 95th-percentile latency (P95), 99th-percentile latency (P99), exponentially weighted moving average (EWMA), support vector classifier (SVC), Google Kubernetes Engine (GKE), recovery point objective (RPO), online transaction processing (OLTP), trusted execution environment (TEE), fifth-generation mobile networking (5G), Internet of Things (IoT), and artificial intelligence (AI). The notation N/A means not applicable.

\begin{table*}[t]
\centering
\caption{Matrix 1: Taxonomy mapping across the 24 included papers. Rows below the rule have a learned decision component.\label{tab:matrix1}}
\resizebox{\textwidth}{!}{%
\begin{tabular}{@{}lccccccc@{}}
\toprule
Paper & Substrate & Objective & Decision mechanism & Impl. layer & Scaling behaviour & Evaluation method & Domain \\
\midrule
\cite{manaouil2021wanwide} & OS container (vanilla) & Pod/API/DNS latency & N/A (stock K8s) & Kernel/interface & N/A & Bare-metal, WAN emulation & Edge, general \\
\cite{kedar2026digitaltwin} & OS container + Kafka & Recovery speed, SLA & Hybrid: quorum + ML ensemble & Streaming/inference/control & Adaptive hybrid & Public-cloud clusters (GCP, AWS) & Manufacturing \\
\cite{tola2026selfhealing} & OS container (K8s/K3s) & SLA risk, MTTR & Predictive loop (LSTM + policy) & User-space middleware & Predictive & Simulation & Telecom/NFV \\
\cite{chen2026empirical} & OS container (K8s) & Response time, success rate & N/A (stock K8s recovery) & Control + data plane & Failure-triggered & Emulation + bare-metal & General cloud/edge \\
\cite{barletta2026slo} & OS container (K8s/Swarm) & Orchestration/failover time & Static, fixed-priority & Control-plane core & Failure-triggered & Bare-metal, two testbeds & Real-time/5G \\
\cite{morabito2026consensus} & OS container (Docker) & Task placement latency & Decentralized consensus & User-space middleware & N/A & Bare-metal (Pi cluster) & Edge microservices \\
\cite{rodriguez2026arapo} & OS container (custom orch.) & Overload time, energy & Hybrid: CNN+LSTM + threshold & Distributed control plane & Predictive & Heterogeneous edge testbed & Healthcare/hospital \\
\cite{gao2026dsfro} & Container + GPU (MIG) & Job completion time & Hybrid ML (LSTM + GA) & Scheduler extender & Predictive+placement & Real testbed + trace replay & Hyperscale ML training \\
\cite{keum2026mutualcloud} & Custom P2P/DHT agent & Scheduling/recovery latency & Decentralized lease + DHT & User-space, custom & N/A & Simulation & General distributed \\
\cite{gurudu2026chaos} & OS container (K8s) & MTTD, MTTR, availability & Static (native K8s healing) & User-space (observability) & Reactive & Case study/demo & General enterprise \\
\cite{chen2026distributed} & OS container (K8s) & P99 latency, consistency & Rule-based, threshold+hysteresis & User-space middleware & N/A & Fault-injection + statistics & E-commerce \\
\cite{bongu2026middleware} & OS container (K8s, dual) & Throughput, p95 latency & Reactive, policy threshold & User-space middleware & Reactive & Bare-metal, dual-cluster & Enterprise hybrid-cloud \\
\cite{kwon2026erpf} & OS container + bare-metal edge & Task duration, std. dev. & Adaptive heuristic (EWMA) & User-space (above K8s) & Partition-level resize & Bare-metal, real failure & Video analytics/edge AI \\
\cite{javens2026adaptive} & OS container (Docker Swarm) & Voting rounds, failure rate & Real-time control loop + SVC & User-space middleware & Adaptive quorum & Emulated Swarm tasks, scripted faults & IoT/edge sensing \\
\cite{kang2026evmm} & OS container (K8s, CRI-O) & Execution time, container restarts, request success & Static priority model & Node-level agent via CRI & Vertical memory scaling, suspend and checkpoint & Bare-metal, four-node cluster & Stateful batch and stateless services \\
\cite{naim2026doublecheck} & Host/device level & Failure freq., downtime & Static, two-step heartbeat & User-space middleware & N/A & 90-day field deployment & Clinical support \\
\cite{jerrari2026ha} & OS container (GKE) & RTO, RPO, throughput & Static (Raft consensus) & User-space middleware & Failure-triggered & Managed cloud (GKE), network emulation & Enterprise OLTP \\
\cite{jin2026replicationfree} & OS container (Alibaba K8s) & RTO, throughput & Static, lease state machine & Kernel + user-space & Failure-triggered & Production cloud cluster & Cloud messaging \\
\cite{li2026pyramid} & Hybrid: VM-TEE + container & Control/data-plane latency & Static, guard state machine & Control plane + kernel & N/A & Bare-metal + simulation & Multi-tenant cloud \\
\cite{chen2026cast} & OS container (service under test) & Vulnerabilities found, bug coverage & Rule-based test prioritisation and threshold oracles & User-space testing framework & N/A & Production deployment and 48-bug benchmark & Production cloud microservices \\
\cite{jensen2021raft} & N/A (network emulation testbed) & Success rate, latency, elections & N/A (studies Raft/etcd) & System coordination substrate & N/A & Network emulation & General coordination \\
\midrule
\cite{dsouza2026selfhealing} & OS container (K8s) & MTTR, success rate, energy & Hybrid: anomaly detection + RL & User-space & Failure-triggered & Fault-injection benchmark & Web applications \\
\cite{nigam2026chaos} & OS container (K8s, Pi, GPU) & SLO violations, MTTR & Deep RL (PPO) & User-space middleware & N/A & Bare-metal, three environments & Mixed: web, edge, AI \\
\cite{mishra2026aidriven} & OS container + OpenStack & Prediction accuracy, MTTR & Hybrid: supervised + threshold & User-space & Predictive & Log analysis + emulated cloud & General hybrid cloud \\
\bottomrule
\end{tabular}}
\end{table*}

Read across substrate and implementation-layer columns, the corpus
concentrates heavily on OS-level containers with user-space
middleware implementation, the same category the implemented mechanism
in Section~VII of the main paper occupies; only three papers,
\cite{barletta2026slo}, \cite{li2026pyramid} and
\cite{jin2026replicationfree}, reach into
control-plane or kernel-level modification, the latter two citing minimal
intrusive change as an explicit design goal, consistent with this
project's own scope boundary.

\medskip\noindent\textit{Matrix 2 -- Empirical Evaluation, Workload, and Data Trace.}\par\smallskip

Matrix~\ref{tab:matrix2} additionally uses Google Cloud Platform (GCP), Amazon Web Services (AWS), round-trip time (RTT), User Datagram Protocol (UDP), Yahoo! Cloud Serving Benchmark (YCSB), input/output (I/O), Advanced Micro Devices (AMD), and Secure Encrypted Virtualization--Secure Nested Paging (SEV-SNP). Product and benchmark names are retained as published by their respective authors.

\begin{table*}[t]
\centering
\caption{Matrix 2: Empirical evaluation methodology and fault-injection parameters. Rows below the rule have a learned decision component.\label{tab:matrix2}}
\resizebox{\textwidth}{!}{%
\begin{tabular}{@{}lcccc@{}}
\toprule
Paper & Environment & Workload generator & Fault-injection tool & Fault parameters \\
\midrule
\cite{manaouil2021wanwide} & Grid'5000, real physical cluster & ClusterLoader + Prometheus & Netem & Latency 0/50/250/400ms; loss up to 10\% \\
\cite{kedar2026digitaltwin} & Real GCP + AWS clusters & Custom workload generator + Kafka & Simulated packet loss/RTT & 150--300ms RTT/1\% loss; 800--1200ms/10\% loss (severe) \\
\cite{tola2026selfhealing} & Simulation (50 synthetic VNFs) & Custom synthetic VNF generator & None & Gradual resource-overload ramp; no network faults \\
\cite{chen2026empirical} & 4/8-node K8s, cloud+edge & Locust & Chaos Mesh & 7 fault types, 4 intensities incl. network partition \\
\cite{barletta2026slo} & 2 heterogeneous testbeds (incl. Pi 4, Kria KV260) & Custom UDP echo, Open5GS+UERANSIM & K8s NoExecution taint & Taint-based exclusion; 5--20 pods/node, 1--3 dest. nodes \\
\cite{morabito2026consensus} & 6$\times$ Raspberry Pi 4 & Custom test harness & None & No fault injection; steady request rates only \\
\cite{rodriguez2026arapo} & 3 heterogeneous edge nodes & Locust & None (load-triggered) & Dual-threshold CPU load only ($\tau_O$=80\%, $\tau_P$=64\%) \\
\cite{gao2026dsfro} & 20-node heterogeneous GPU testbed & Trace replay & None (designed, untested) & Heartbeat 1s, 3-miss threshold; not experimentally triggered \\
\cite{keum2026mutualcloud} & Simulation, single dev machine & Own simulator & Simulated churn model & 1--5\% node failure rate; K8s/Hyperledger baselines modeled, not run \\
\cite{gurudu2026chaos} & K8s (cluster details not given) & Own telemetry & LitmusChaos & Named, not quantified (pod term., CPU/mem., network) \\
\cite{chen2026distributed} & K8s 1.29 + Kafka + PostgreSQL & k6 & Chaos Mesh & 5 fault types incl. partition; duration not quantified \\
\cite{bongu2026middleware} & Private + public K8s clusters & DeathStarBench, Train-Ticket & Not named & Broker outages, cross-cluster link failures \\
\cite{kwon2026erpf} & K8s + real Pi/Jetson devices & Real workloads (YOLOv8, transcoding) & Custom worker-disconnect injector & Real failure across 3 worker pairs, 20s reintegration \\
\cite{javens2026adaptive} & 70 emulated Swarm tasks; hardware not reported & Custom query client & Custom probabilistic injector & 1\% message-drop; 0--32/70 nodes faulty \\
\cite{kang2026evmm} & Four-node bare-metal K8s, 1~Gbps & Stateful benchmark applications; concurrent request workers & None & Memory pressure induced by workload: oversubscription up to 150\%; normal, peak, and overload concurrency \\
\cite{naim2026doublecheck} & 2 real production servers & Live clinical traffic & Manual (by hand) & 4 fault types, manual trigger \\
\cite{jerrari2026ha} & 3-node GKE & pgbench & tc netem & Pod kill; 100\% loss; kill 2/3 nodes in 5s \\
\cite{jin2026replicationfree} & Real Alibaba Cloud K8s & Custom RocketMQ generator & ChaosBlade + manual port stop & Node crash, forced migration, I/O hang, port-stop ``partition'' \\
\cite{li2026pyramid} & 6-node AMD SEV-SNP cluster & K-Bench, app benchmarks & None (threat model only) & Arbitrary disconnection/clock manipulation, not fault-injected \\
\cite{chen2026cast} & Huawei Cloud production services; per-test isolated containers & Replayed production traffic & Application-level fault library & Software exceptions and inter-service communication errors; 137 potential vulnerabilities, 89 confirmed; 90\% coverage of 48 reproduced bugs \\
\cite{jensen2021raft} & Mininet, real etcd binary & Custom write-traffic generator & iptables & Leader crash, full/partial/intermittent partition \\
\midrule
\cite{dsouza2026selfhealing} & Infrastructure not reported & Custom Python monitoring service & Custom (unnamed) & 6 fault types incl. network delay; $\geq$30 trials each \\
\cite{nigam2026chaos} & 6-node K8s, 4-node Pi 5 edge, 4-node A100 & Locust, wrk2 & RL-driven fault selection & Latency 50--300ms, loss 0.1--5\%, CPU 60--95\% \\
\cite{mishra2026aidriven} & Log analysis + emulated hybrid cloud & Not specified & Not specified & No fault parameters reported \\
\bottomrule
\end{tabular}}
\end{table*}

The fault-parameters column is the clearest evidence for the gap
formalised in Section~VII of the main paper: of the 24 coded papers, only
\cite{chen2026empirical}, \cite{chen2026distributed}, \cite{jerrari2026ha},
\cite{jin2026replicationfree}, and \cite{jensen2021raft} inject
anything resembling a genuine network partition, and of these, only
\cite{jerrari2026ha} and \cite{jensen2021raft} separately compare
partition against crash on identical infrastructure. None of these partition-injection studies evaluates a general-purpose Kubernetes decision stage that combines post-suspicion reachability evidence, workload characteristics, and recovery safety before authorising recovery. This narrower statement is intentional because other papers in the corpus do implement secondary verification or safeguards in different contexts.
\medskip\noindent\textit{Matrix 3 -- Causality Gap and Systems Evaluation Audit.}\par\smallskip

\begin{table*}[!t]
\centering
\caption{Matrix 3: Evaluation-quality appraisal. Y = fell into the trap (0); P = partial (0.5); N = avoided (1); N/A = not applicable and not scored. Score = total over applicable criteria; the reason is given for every Y or P cell. Rows below the first rule have a learned decision component; the last row applies the same criteria to this work.\label{tab:matrix3}}
\footnotesize
\setlength{\tabcolsep}{3pt}
\begin{tabular}{@{}>{\RaggedRight\arraybackslash}p{0.07\textwidth}>{\centering\arraybackslash}p{0.07\textwidth}>{\centering\arraybackslash}p{0.07\textwidth}>{\centering\arraybackslash}p{0.08\textwidth}>{\centering\arraybackslash}p{0.05\textwidth}>{\RaggedRight\arraybackslash}p{0.56\textwidth}@{}}
\toprule
Paper & Averages fallacy & Localhost fallacy & Infinite bandwidth & Score & Reason for Y or P \\
\midrule
\cite{manaouil2021wanwide} & N & N & N & 3/3 & -- \\
\cite{kedar2026digitaltwin} & N & N & N & 3/3 & -- \\
\cite{tola2026selfhealing} & N & N/A & N/A & 1/1 & -- \\
\cite{chen2026empirical} & N & N & N & 3/3 & -- \\
\cite{barletta2026slo} & N & N & N/A & 2/2 & -- \\
\cite{morabito2026consensus} & Y & N & N/A & 1/2 & Placement latency reported without variance. \\
\cite{rodriguez2026arapo} & N & N & P & 2.5/3 & Links between nodes neither throttled nor characterised. \\
\cite{gao2026dsfro} & P & N & N & 2.5/3 & Median and p99 for scheduling delay only; other metrics are single averages. \\
\cite{keum2026mutualcloud} & P & Y & N/A & 0.5/2 & p95 without repeat count or variance across seeds. All results simulated on one machine. \\
\cite{gurudu2026chaos} & Y & N/A & N/A & 0/1 & Single values without repetition or spread. \\
\cite{chen2026distributed} & N & N/A & N/A & 1/1 & -- \\
\cite{bongu2026middleware} & P & N & N/A & 1.5/2 & p95 latency only; throughput, availability and loss are single values. \\
\cite{kwon2026erpf} & P & N & N & 2.5/3 & Recovery timings averaged without stated trial counts. \\
\cite{javens2026adaptive} & N & P & N/A & 1.5/2 & Message loss is a modelled probability, not a degraded path. \\
\cite{kang2026evmm} & N & N & N & 3/3 & -- \\
\cite{naim2026doublecheck} & Y & N & N & 2/3 & Average detection times over five trials without spread. \\
\cite{jerrari2026ha} & N & N & N & 3/3 & -- \\
\cite{jin2026replicationfree} & Y & N & N/A & 1/2 & Single RTO values without stated trial counts. \\
\cite{li2026pyramid} & Y & N & N/A & 1/2 & Single values per tenant count; no repeats or error bars. \\
\cite{chen2026cast} & N & N & N/A & 2/2 & -- \\
\cite{jensen2021raft} & P & N & N & 2.5/3 & Fault traces show one of five runs without variance. \\
\midrule
\cite{dsouza2026selfhealing} & P & N/A & N/A & 0.5/1 & Medians with 95\% CI for two of six fault types only. \\
\cite{nigam2026chaos} & N & N & N & 3/3 & -- \\
\cite{mishra2026aidriven} & Y & N/A & N/A & 0/1 & Single point values; baselines not described. \\
\midrule
This work & N & N & N & 3/3 & Per-run values, SD, 95\% CI and range reported; separate virtual machines over real interfaces; degradation applied with \texttt{tc}. \\
\bottomrule
\end{tabular}
\end{table*}

The averages fallacy is the most frequently encountered trap, present or partial in 12 of the 24 papers; the
localhost fallacy is comparatively rare, present cleanly in only
\cite{keum2026mutualcloud}, whose evaluation is a full simulation
with no real network path at all. Of the two papers most central
to the gap identified in this review, \cite{jerrari2026ha} avoids every applicable trap and \cite{jensen2021raft} is partial on the averages fallacy
only.

\medskip\noindent\textit{Matrix 4 -- Sovereignty and Resource-Constrained Settings Fit.}\par\smallskip

\begin{table}[t]
\centering
\caption{Matrix 4: Sovereignty and resource-constrained fit (1--5 scale; anchors in Table~\ref{tab:matrix4-rubric}). Rows are ordered by the sum of the four sub-scores.\label{tab:matrix4}}
\resizebox{\columnwidth}{!}{%
\begin{tabular}{@{}lcccc@{}}
\toprule
Paper & Offline autonomy & Bandwidth/egress & Data residency & Hardware efficiency \\
\midrule
\cite{rodriguez2026arapo} & 4 & 4 & 5 & 4 \\
\cite{keum2026mutualcloud} & 4 & 2 & 3 & 3 \\
\cite{kwon2026erpf} & 3 & 3 & 1 & 5 \\
\cite{naim2026doublecheck} & 3 & 3 & 1 & 4 \\
\cite{manaouil2021wanwide} & 1 & 4 & 3 & 2 \\
\cite{kedar2026digitaltwin} & 2 & 3 & 2 & 3 \\
\cite{chen2026empirical} & 2 & 4 & 1 & 3 \\
\cite{javens2026adaptive} & 2 & 3 & 1 & 4 \\
\cite{jensen2021raft} & 2 & 3 & 1 & 3 \\
\cite{barletta2026slo} & 2 & 2 & 1 & 4 \\
\cite{bongu2026middleware} & 1 & 2 & 3 & 2 \\
\cite{nigam2026chaos} & 2 & 3 & 1 & 2 \\
\cite{li2026pyramid} & 1 & 1 & 4 & 1 \\
\cite{morabito2026consensus} & 1 & 2 & 1 & 3 \\
\cite{tola2026selfhealing} & 2 & 1 & 1 & 3 \\
\cite{jerrari2026ha} & 1 & 2 & 1 & 2 \\
\cite{kang2026evmm} & 1 & 1 & 1 & 3 \\
\cite{jin2026replicationfree} & 1 & 1 & 1 & 2 \\
\cite{dsouza2026selfhealing} & 1 & 1 & 1 & 2 \\
\cite{chen2026cast} & 1 & 1 & 1 & 1 \\
\cite{gao2026dsfro} & 1 & 1 & 1 & 1 \\
\cite{gurudu2026chaos} & 1 & 1 & 1 & 1 \\
\cite{chen2026distributed} & 1 & 1 & 1 & 1 \\
\cite{mishra2026aidriven} & 1 & 1 & 1 & 1 \\
\bottomrule
\end{tabular}}
\end{table}

\begin{table}[t]
\centering
\caption{Scoring anchors for Matrix~4. Scores of 2 and 4 lie between the neighbouring anchors.\label{tab:matrix4-rubric}}
\footnotesize
\setlength{\tabcolsep}{3pt}
\begin{tabular}{@{}>{\RaggedRight\arraybackslash}p{0.21\columnwidth}>{\RaggedRight\arraybackslash}p{0.235\columnwidth}>{\RaggedRight\arraybackslash}p{0.235\columnwidth}>{\RaggedRight\arraybackslash}p{0.235\columnwidth}@{}}
\toprule
Criterion & 1 & 3 & 5 \\
\midrule
Offline autonomy & No offline or disconnected operation considered & Local operation designed for or tested under short disruptions & Sustained disconnected operation designed for and demonstrated \\
Bandwidth or egress & Not considered & Constrained or high-latency links part of the evaluation & Bandwidth or egress reduction a primary, measured objective \\
Data residency & Not discussed & Data locality or routing addressed indirectly & Local residency of sensitive data required by design \\
Hardware efficiency & Data-centre hardware; no constraint considered & Edge deployability profiled, or modest commodity hardware & Constrained devices used, with overhead measured \\
\bottomrule
\end{tabular}
\end{table}

\cite{rodriguez2026arapo} is the only paper scoring 4 or above on all four
sub-dimensions, reflecting its explicit rural-care,
intermittent-connectivity motivation. No paper scores above 3 on
data residency except \cite{rodriguez2026arapo} and
\cite{li2026pyramid}, the latter for reasons orthogonal to this
project's concerns, namely protecting tenant data from an untrusted
cloud provider via hardware trusted execution, rather than supporting offline clinical operation. This gap between
high-hardware-efficiency papers (\cite{kwon2026erpf}, scoring 5 on that sub-dimension) and
high-residency papers is itself notable: the reviewed literature does
not present a single mechanism strong on both constrained-hardware
efficiency and data residency simultaneously, a combination directly
relevant to the healthcare deployment context of this project.




\clearpage
\section{Run-Level OpenMRS Tables}
These tables give the run-level OpenMRS values summarised in Section~IX of the main paper.

\begin{table}[h]
\centering
\caption{OpenMRS Scenario B: original backend Pod identity in the 15 per-configuration runs.\label{tab:openmrs-preservation}}
\begin{tabular}{@{}lrrr@{}}
\toprule Mode & Runs & Pod changed & Pod retained\\ \midrule
V0 & 3 & 3 & 0\\
T5 & 3 & 3 & 0\\
T10 & 3 & 3 & 0\\
T30 & 3 & 3 & 0\\
FA1 (initial) & 3 & 2 & 1\\ \bottomrule
\end{tabular}
\end{table}

\begin{table}[h]
\centering
\caption{OpenMRS Scenario B Failure-Aware runs: decision chain from verification to Pod outcome. Fault durations are from each run's recorded fault timestamps.\label{tab:openmrs-decision-chain}}
\footnotesize
\begin{tabular}{@{}>{\RaggedRight\arraybackslash}p{0.09\textwidth}>{\RaggedRight\arraybackslash}p{0.12\textwidth}>{\RaggedRight\arraybackslash}p{0.10\textwidth}>{\RaggedRight\arraybackslash}p{0.09\textwidth}>{\RaggedRight\arraybackslash}p{0.13\textwidth}>{\RaggedRight\arraybackslash}p{0.245\textwidth}>{\RaggedRight\arraybackslash}p{0.10\textwidth}@{}}
\toprule
Run & Verification & Decision & Toleration & Fault (s) & Native lifecycle events & Original Pod\\ \midrule
003 & 3/3 HTTP 200 & PRESERVE & 300~s & 121.2 & Taint-manager deletion cancelled & Retained\\
004 & 3/3 HTTP 200 & PRESERVE & 300~s & 902.4 & Taint-manager deletion; ReplicaSet created replacement on Worker2 & Replaced\\
005 & 3/3 HTTP 200 & PRESERVE & 300~s & 468.7 & Taint-manager deletion; ReplicaSet created replacement on Worker2 & Replaced\\
006 (validation) & 3/3 HTTP 200 & PRESERVE & Indefinite & 468.0 & None disruptive; Pod Running and Ready on Worker1 after the run & Retained\\ \bottomrule
\end{tabular}
\end{table}

\begin{table}[h]
\centering
\caption{OpenMRS component resource audit; usage is one observed snapshot.\label{tab:openmrs-resources}}
\footnotesize
\setlength{\tabcolsep}{2.5pt}
\begin{tabular}{@{}lrrrr@{}}
\toprule Component & Observed & Observed & CPU & Memory\\
 & CPU (m) & memory (MiB) & request (m) & request (MiB)\\ \midrule
Backend & 19 & 929 & 750 & 1400\\
MariaDB & 9 & 130 & 300 & 512\\
Frontend & 1 & 2 & 100 & 96\\
Gateway & 1 & 3 & 100 & 96\\ \midrule
Total & 30 & 1064 & 1250 & 2104\\ \bottomrule
\end{tabular}
\end{table}

\begin{table}[h]
\centering
\caption{Hypothetical additional OpenMRS backend memory reservation by worker (request-based snapshot, not observed concurrent duplication).\label{tab:openmrs-worker-capacity}}
\footnotesize
\begin{tabular}{@{}lrrrl@{}}
\toprule Worker & Allocatable & Existing & With backend & Fits\\
 & (MiB) & requests (MiB) & (MiB) & \\ \midrule
Worker1 & 3915 & 2992 & 4392 & No\\
Worker2 & 3915 & 1024 & 2424 & Yes\\
Worker3 & 1968 & 0 & 1400 & Yes\\
Worker4 & 1968 & 704 & 2104 & No\\
\bottomrule
\end{tabular}
\end{table}

\begin{table}[h]
\centering
\caption{OpenMRS Scenario B clinical probe failures in the nine runs with samples. The interval runs from the first failed probe to the next successful one; ``Open'' means no successful probe followed, so the interval is censored.\label{tab:openmrs-clinical-episodes}}
\footnotesize
\begin{tabular}{@{}lrrl@{}}
\toprule
Run & Failed probes & Interval (s) & Original Pod\\ \midrule
V0/003 & 9 & 11.847 & Replaced\\
V0/004 & 14 & Open & Replaced\\
T5/003 & 6 & Open & Replaced\\
T5/004 & 17 & Open & Replaced\\
T10/003 & 6 & 17.795 & Replaced\\
T10/004 & 24 & Open & Replaced\\
T30/003 & 14 & 16.965 & Replaced\\
FA1/003 & 1 of 102 & 1.209 & Retained\\
FA1/004 & 19 of 677 & 23.269 & Replaced\\ \bottomrule
\end{tabular}
\end{table}

\begin{table}[h]
\centering
\caption{OpenMRS Scenario C: fault duration, time from fault to first clinical restoration, and time from fault clearance to restoration, in seconds. Descriptive only.\label{tab:openmrs-scenario-c}}
\footnotesize
\setlength{\tabcolsep}{4pt}
\begin{tabular}{@{}lrrr@{}}
\toprule
Run & Fault & Fault to & Clearance to\\
 & duration & restoration & restoration\\ \midrule
V0/003 & 902.410 & 922.216 & 19.807\\
V0/004 & 908.435 & 927.227 & 18.792\\
T5/003 & 910.445 & 935.102 & 24.657\\
T5/004 & 904.425 & 931.693 & 27.268\\
T10/003 & 132.502 & 592.121 & 459.618$^{a}$\\
T10/004 & 903.221 & 919.773 & 16.551\\
T30/003 & 908.235 & 931.060 & 22.825\\
T30/004 & 905.172 & 927.932 & 22.761\\
FA1/003 & 609.550 & 635.936 & 26.386\\
FA1/004 & 903.688 & 922.523 & 18.835\\ \midrule
\multicolumn{4}{@{}l}{Configuration means ($n=2$)}\\
V0 & 905.423 & 924.722 & 19.300\\
T5 & 907.435 & 933.398 & 25.963\\
T10 & 517.862 & 755.947 & 238.085\\
T30 & 906.704 & 929.496 & 22.793\\
FA1 & 756.619 & 779.230 & 22.611\\ \bottomrule
\multicolumn{4}{@{}l}{\scriptsize $^{a}$Cause of the interval not verified.}
\end{tabular}
\end{table}

\begin{table}[h]
\centering
\caption{OpenMRS Scenario C+D intervals (s), primary cohort $n=9$ (runs 011 and 013--020).\label{tab:openmrs-cd}}
\footnotesize
\setlength{\tabcolsep}{2pt}
\begin{tabular}{@{}lccc@{}}
\toprule
Interval & Mean $\pm$ SD & Median & Range \\
\midrule
Fault $\rightarrow$ database probe unreachable$^{\dagger}$ & 3.98 $\pm$ 0.79 & 3.53 & 3.45--5.83 \\
Fault $\rightarrow$ RECOVER decision & 57.57 $\pm$ 2.29 & 57.11 & 54.30--61.45 \\
Decision $\rightarrow$ fence verified & 1.80 $\pm$ 0.24 & 1.74 & 1.63--2.43 \\
Fence verified $\rightarrow$ recovery verified & 102.10 $\pm$ 8.74 & 106.23 & 89.41--111.37 \\
Fault $\rightarrow$ recovery verified & 161.48 $\pm$ 8.95 & 165.16 & 149.53--170.95 \\
Reconciliation & 575.56 $\pm$ 26.60 & 584.00 & 505--588 \\
\bottomrule
\multicolumn{4}{@{}p{0.9\textwidth}@{}}{\scriptsize $^{\dagger}$The independent database probe on the verification path, which loses the database within seconds of the fault. Native suspicion (\texttt{Ready=Unknown}) follows about 50~s later, and only then does the controller verify.}\\
\end{tabular}
\end{table}


\FloatBarrier

\section{Experiment Inventory and Second-Phase Campaigns}
\setcounter{table}{12}
\begin{table}[!t]
\centering
\caption{Experiment inventory for both evaluations: attempted, aborted and analysed runs. Second-phase runs are listed in Section~S3.}\label{tab:s13}
\footnotesize
\setlength{\tabcolsep}{3pt}
\begin{tabular}{@{}>{\RaggedRight\arraybackslash}p{0.09\textwidth}>{\RaggedRight\arraybackslash}p{0.17\textwidth}>{\RaggedRight\arraybackslash}p{0.20\textwidth}>{\RaggedRight\arraybackslash}p{0.09\textwidth}>{\RaggedRight\arraybackslash}p{0.38\textwidth}@{}}
\toprule
Workload & Campaign & Configurations & Runs & Note \\
\midrule
Nginx & Scenario B & Native, FA1 & 30, 32 & Control-plane partition; workload reachable. FA1 run 019 aborted after 3.8~s, before suspicion. \\
Nginx & Scenario B demonstration & FA1 & 4 & Runs 201, 301, 302, 501: all retained with PRESERVE; not analysed. \\
Nginx & Scenario C & Native, FA1 & 30, 30 & Control-plane and client paths isolated. \\
Nginx & Scenario A & Native, FA1 & 5, 5 & Worker1 shut down; 6 and 7 attempts, incomplete attempts excluded. \\
Nginx & Single-node baseline & No orchestration & 10 & Host Nginx on Worker1, observed in every Scenario A run. \\
Nginx & Scenario D (PostgreSQL) & FA1 & 31 & Old writer stopped before fencing in every run. \\
Nginx & Scenario E & FA1: mild, moderate, severe & 10, 10, 3 & Control-plane path degraded on the worker's interface. \\
Nginx & Overhead & Native, FA1 & 30, 30 & Normal operation; first-phase campaign. \\
Nginx & Candidate selection & FA1 & 4 & Controlled cases with configured reliability scores. \\
OpenMRS & Scenario B & V0, T5, T10, T30, FA1 & 15 & Three per configuration. \\
OpenMRS & Scenario B run 006 & FA1 & 1 & Validation with non-expiring tolerations, reported separately. \\
OpenMRS & Scenario C & V0, T5, T10, T30, FA1 & 10 & Two per configuration. \\
OpenMRS & Scenario E & Five configurations, three profiles & 14 & Severe observations for V0, T5, T30, and FA1. \\
OpenMRS & Scenario C+D runs 011--020 & FA1 & 10 & MariaDB on NFS with storage fence and SSH reconciliation; run 012 interrupted after recovery. \\
OpenMRS & Scenario C+D run 010 & FA1 & 0 & Stopped at preflight (Worker3 cordoned); no fault applied. \\
OpenMRS & Scenario C+D pilots & FA1 & 4 & Fencing and recovery verified in all four; not pooled. \\
OpenMRS & Scenario C+D run 003 & FA1 & 1 & Trial run before the campaign, without the reconciliation stage; the campaign runs are numbered from 010. \\
\bottomrule
\end{tabular}
\end{table}

\begin{table}[ht]
\centering
\caption{Second-phase campaigns (commit \texttt{4ce8078}, extension addendum \texttt{ddfc677}).}\label{tab:s14}
\footnotesize
\begin{tabular}{@{}p{0.22\textwidth}p{0.30\textwidth}p{0.10\textwidth}p{0.30\textwidth}@{}}
\toprule
Campaign & Arms and conditions & Runs & Status \\
\midrule
Baseline ladder (seed 3106524066) & V0, T10, NX, R2, FA1; reachable 120 and 600~s, unreachable 480~s & 75 & 75 completed, 0 aborted \\
Ladder, second set (seed 781441089) & V0, T10, NX, R2, FA1; reachable 600~s, unreachable 480~s & 30 & 30 completed, 0 aborted; pooled with the baseline ladder \\
Ladder, third set (seed 472857209) & V0, T10, NX, R2, FA1; reachable 120~s & 15 & 15 completed, 0 aborted; pooled with the baseline ladder \\
Probe grid & Six workload states, 100 trials each, $k=1,3,5$ & 600 trials & Completed \\
Overhead & 30 idle windows with and 30 without the controller & 60 windows & Completed \\
Live writer, first campaign (seed 3063071617) & Planned 10 fenced and 10 unfenced & 20 & 20 completed; no fence applied in any run; reported as unfenced \\
Live writer, fenced (seed 721711068) & Fenced & 10 & 10 completed, 0 aborted \\
Live writer, fenced (seed 2559281726) & Fenced & 10 & 10 completed, 0 aborted \\
Probe grid, extended (seed 781441089) & Reply delay 0.5, 1.0, 1.5~s; reply loss 70 and 90\%; 100 trials each, $k=1,3,5$ & 500 trials & Completed \\
Extension E1 (seed 2000988837) & T5, T30; reachable 120 and 600~s & 20 & 20 completed, 0 aborted \\
Extension E2 (seed 3566506833) & V0, T10, NX, FA1; reachable 15, 60, 300~s & 36 & 36 completed, 0 aborted \\
OpenMRS Scenario B under load (runs 101--103) & V0, T5, T10, T30, FA1; 600-s partition; ten virtual clinical users & 14 & 14 completed, 0 aborted (T30: two runs) \\
\bottomrule
\end{tabular}
\end{table}


\section{Original-Campaign Result Tables}
The tables in this section are referred to from the main text by number.

\begin{table}[ht]
\centering
\caption{Scenario B aggregate results. Medians carry BCa bootstrap 95\% intervals (10,000 resamples); p95 and p99 are given where $n\ge20$; with $n\approx30$, p99 lies near the maximum. A dash denotes an inapplicable or unavailable value. FA1 run 019 (aborted before suspicion) is excluded. $^{\S}$Run 003 is excluded: no observation was recorded for 625.9~s, so its outage is unobserved.\label{tab:scenario-b-results}}
\resizebox{\linewidth}{!}{%
\begin{tabular}{@{}llrrrrrrr@{}}
\toprule
Mode & Metric & $n$ & Mean (s) & SD (s) & Median [BCa 95\% CI] (s) & p95 (s) & p99 (s) & Range (s) \\
\midrule
Native & Replacement pending & 24 & 347.83 & 3.28 & 347.24 [345.64, 349.31] & 353.28 & 353.56 & [341.66, 353.58] \\
Native & Replacement running & 30 & 349.87 & 3.62 & 349.76 [347.77, 352.02] & 355.03 & 355.79 & [342.83, 355.91] \\
Native & HTTP outage & 8 & 4.57 & 0.37 & 4.42 [4.29, 4.87] & -- & -- & [4.22, 5.34] \\
Failure-Aware & Suspicion delay & 32 & 51.36 & 18.72 & 47.52 [46.94, 49.33] & 54.45 & 122.18 & [41.76, 152.29] \\
Failure-Aware & Decision delay & 32 & 20.94 & 110.37 & 1.35 [1.20, 1.42] & 2.23 & 432.52 & [1.12, 625.79] \\
Failure-Aware & HTTP outage$^{\S}$ & 4 & 115.71 & 90.23 & 114.55 [6.57, 227.19] & -- & -- & [6.57, 227.19] \\
Failure-Aware & Replacement/service recovery & 0 & -- & -- & -- & -- & -- & -- \\
\bottomrule
\end{tabular}}
\end{table}

\begin{table}[ht]
\centering
\caption{Scenario C aggregate results. Medians carry BCa bootstrap 95\% intervals (10,000 resamples); p95 and p99 are given where $n\ge20$; with $n=30$, p99 lies near the maximum.\label{tab:scenario-c-results}}
\resizebox{\linewidth}{!}{%
\begin{tabular}{@{}llrrrrrrr@{}}
\toprule
Mode & Metric & $n$ & Mean (s) & SD (s) & Median [BCa 95\% CI] (s) & p95 (s) & p99 (s) & Range (s) \\
\midrule
Native & Replacement pending & 25 & 349.14 & 3.69 & 349.43 [346.54, 350.54] & 355.71 & 355.95 & [342.98, 356.03] \\
Native & Replacement running & 30 & 353.97 & 18.49 & 350.65 [347.89, 352.38] & 358.26 & 423.25 & [344.17, 449.77] \\
Native & End-to-end recovery (ERL) & 30 & 356.56 & 18.36 & 352.84 [350.80, 355.63] & 361.80 & 425.47 & [346.11, 451.22] \\
Native & HTTP outage (SID) & 30 & 374.66 & 116.70 & 353.10 [351.31, 355.44] & 361.64 & 809.62 & [347.24, 992.13] \\
Failure-Aware & Suspicion delay (FDL) & 30 & 50.15 & 3.78 & 51.47 [48.58, 52.19] & 55.10 & 55.61 & [43.05, 55.70] \\
Failure-Aware & Decision delay & 30 & 10.81 & 1.06 & 10.26 [10.13, 10.34] & 12.85 & 13.18 & [10.08, 13.28] \\
Failure-Aware & Recovery-release delay & 30 & 0.08 & 0.16 & 0.03 [0.02, 0.06] & 0.21 & 0.69 & [0.01, 0.86] \\
Failure-Aware & Replacement pending & 18 & 61.80 & 3.31 & 62.51 [59.51, 63.59] & -- & -- & [55.05, 67.01] \\
Failure-Aware & Replacement running & 30 & 63.03 & 3.46 & 64.12 [62.35, 64.74] & 67.31 & 68.02 & [55.03, 68.31] \\
Failure-Aware & End-to-end recovery (ERL) & 30 & 65.43 & 3.22 & 66.11 [65.12, 66.63] & 69.22 & 71.92 & [57.66, 73.01] \\
Failure-Aware & HTTP outage (SID) & 30 & 62.69 & 3.66 & 62.99 [60.88, 64.08] & 68.38 & 68.78 & [55.50, 68.90] \\
\bottomrule
\end{tabular}}
\end{table}

\begin{table}[ht]
\centering
\caption{Scenario D storage-authority correctness and timing results. Intervals are Student-$t$ 95\% intervals for the mean.\label{tab:scenario-d-results}}
\resizebox{\linewidth}{!}{%
\begin{tabular}{@{}lrrrrr@{}}
\toprule
Metric & $n$ & Mean (s) & SD (s) & 95\% CI (s) & Range (s) \\
\midrule
Old-authority stopping time & 31 & 3.38 & 0.78 & [3.09, 3.67] & [2.05, 4.69] \\
Fence-to-recovery-release time & 31 & 11.88 & 9.79 & [8.29, 15.47] & [9.09, 64.36] \\
Recovery-release-to-service time & 31 & 3.93 & 1.22 & [3.48, 4.38] & [2.46, 9.03] \\
Fence-to-service time & 31 & 15.82 & 10.82 & [11.85, 19.79] & [11.67, 73.39] \\
\bottomrule
\end{tabular}}
\end{table}

\begin{table}[ht]
\centering
\caption{Controlled multi-candidate recovery-selection checks. Each case is a single controlled run on the live cluster with supplied reliability inputs.\label{tab:candidate-selection-results}}
\resizebox{\linewidth}{!}{%
\begin{tabular}{@{}p{3.0cm}p{5.0cm}p{3.2cm}p{5.2cm}@{}}
\toprule
Test & Controlled condition & Selected candidate & Observed outcome \\
\midrule
Reliability selection & All candidates Ready, schedulable, and resource-feasible & \texttt{k3s-worker3} & Highest reliability score (0.997000) selected despite \texttt{k3s-worker2} exposing twice its allocatable CPU. \\
Resource exclusion & \texttt{k3s-worker3} requested 1950m of 2000m CPU, leaving 50m for a 100m recovery requirement & \texttt{k3s-worker4} & Worker3 was correctly marked ineligible despite having the highest reliability score. \\
Availability exclusion & \texttt{k3s-worker3} remained Ready but was cordoned (\texttt{schedulable=false}) & \texttt{k3s-worker4} & The highest-reliability node was excluded and the next-best eligible candidate was selected. \\
RECOVER placement & All candidates eligible; original workload classified \texttt{Unreachable} after three failed probes & \texttt{k3s-worker3} & Controller emitted RECOVER, constrained the Deployment to Worker3, released the Worker1 Pod, and the replacement became Running on Worker3. \\
\bottomrule
\end{tabular}}
\end{table}

\begin{table}[ht]
\centering
\caption{Scenario E outcomes under degraded control-plane networking.\label{tab:scenario-e-results}}
\resizebox{\linewidth}{!}{%
\begin{tabular}{@{}lrrrrrr@{}}
\toprule
Profile & Runs & \texttt{Ready=Unknown} & PRESERVE & Replacement & Original Pod retained & Decision scored \\
\midrule
Mild & 10 & 0/10 & Not triggered & 0/10 & 10/10 & Not scored (non-event) \\
Moderate & 10 & 0/10 & Not triggered & 0/10 & 10/10 & Not scored (non-event) \\
Severe & 3 & 3/3 & 3/3 & 0/3 & 3/3 & 3/3 correct \\
\bottomrule
\end{tabular}}
\end{table}

\begin{table}[ht]
\centering
\caption{OpenMRS Scenario B fault duration of each run, from the recorded fault-start and fault-clear timestamps, against the approximate native eviction window.\label{tab:openmrs-fault-window}}
\footnotesize
\setlength{\tabcolsep}{3pt}
\begin{tabular}{@{}lrrrrr@{}}
\toprule
Mode & Window (s) & Run 003 & Run 004 & Run 005 & Retained\\ \midrule
V0 & $\approx$350 & 902.5 & 422.3 & 460.1 & 0/3\\
T5 & $\approx$55 & 122.3 & 422.4 & 468.4 & 0/3\\
T10 & $\approx$60 & 602.8 & 422.7 & 459.8 & 0/3\\
T30 & $\approx$80 & 902.3 & 460.1 & 467.5 & 0/3\\
FA1 (initial) & $\approx$350 & 121.2 & 902.4 & 468.7 & 1/3 (run 003)\\ \bottomrule
\end{tabular}
\end{table}

\begin{table}[ht]
\centering
\caption{Principal findings with exact (Clopper--Pearson) 95\% intervals. ``Window'' is the native grace period plus the Pod's toleration.\label{tab:principal-findings}}
\footnotesize
\begin{tabular}{@{}>{\RaggedRight\arraybackslash}p{0.04\linewidth}>{\RaggedRight\arraybackslash}p{0.29\linewidth}>{\RaggedRight\arraybackslash}p{0.32\linewidth}>{\RaggedRight\arraybackslash}p{0.27\linewidth}@{}}
\toprule
& Finding & Evidence & Strength\\ \midrule
F1 & Native Kubernetes replaces a reachable workload once a partition outlasts the window. & Nginx: 30/30 replaced (retention 0\%, 95\% CI 0--11.6\%). OpenMRS: 12/12 native and tuned runs replaced (0--26.5\%); 14/14 runs longer than their window replaced. & Strong; repeated on two workloads.\\
F2 & The controller classifies a reachable workload correctly and chooses PRESERVE. & Nginx: PRESERVE in 32 of 32 valid runs (95\% CI 89.1--100\%), no RECOVER; one further attempt (run 019) was aborted before suspicion could occur. OpenMRS: 12/12 verification attempts succeeded; PRESERVE in 4/4 runs (95\% CI 39.8--100\%). & Strong for Nginx; consistent but small for OpenMRS.\\
F3 & The integrated configuration retains the workload. & Nginx: 32/32 retained (89.1--100\%) against 0/30 natively, a difference of 100 points (84.4--100), Fisher exact $p=2.2\times10^{-18}$. OpenMRS run 006 repeats the Nginx outcome on the healthcare workload. & Strong; established on Nginx and replicated on OpenMRS.\\
F4 & A PRESERVE decision does not by itself stop native eviction. & OpenMRS: both FA1 runs that outlasted a 300-second window were evicted despite PRESERVE. & Clear mechanism; two runs.\\
F5 & Avoiding a duplicate backend avoids a substantial reservation. & OpenMRS backend: 929~MiB used, 1,400~MiB and 750m CPU requested. Every replacement (12 native and tuned runs, FA1 runs 004 and 005) was placed on Worker2, taking 48\% of its unreserved memory; Worker4 could not admit one. & Request-based; client cost under load in Scenario B runs 101--103.\\
F6 & The controller releases an unreachable workload for recovery on an eligible candidate. & Nginx: RECOVER and a running replacement in 30/30 runs. OpenMRS: RECOVER in 2/2 runs, replacement placed on the selected worker, clinical service restored. & Strong for Nginx; consistent for OpenMRS.\\
F7 & Storage is fenced before a duplication-sensitive workload is released and stays fenced until the original writer is retired. & PostgreSQL: ordering held in 31/31 runs (88.8--100\%) with the writer stopped first. MariaDB: fence verified before replacement and released only after the original process was verified absent in 10/10 runs (69.2--100\%), with the original runtime not stopped first. & Ordering held in every run.\\
F8 & Removal of a Pod from Kubernetes does not mean its process has exited. & MariaDB run 003: original process in state \texttt{D} for more than 14 hours after its Pod disappeared from the API; runs 011--020 retired it by SSH reboot before releasing the fence (mean reconciliation 575.56~s). & Observed; reconciliation repeatable in 10 runs.\\
F9 & Under genuine worker shutdown the controller classifies the workload correctly and, when a verified destination exists, restores service sooner than native K3s or a single node. & Nginx Scenario A: unreachable in 5/5; RECOVER with replacement on the selected node in 4/5 (28.4--99.5\%), median outage 61.36~s over those four, against 350.06~s natively and 530.88~s for a single node. In run 004 no verified destination was available, the workload was held, and the outage was 592.37~s; the five-run median is 62.43~s. & Five runs per arm; wide intervals.\\ \bottomrule
\end{tabular}
\end{table}

\input{bibliography_block}
\end{document}
